% Part: history
% Chapter: biographies 
% Section: alonzo-church

\documentclass[../../../include/open-logic-section]{subfiles}

\begin{document}

\olfileid[hi]{his}{bio}{chu}

\olsection{अलॉन्ज़ो चर्च}

\olphoto{church-alonzo}{अलॉन्ज़ो चर्च}

अलॉन्ज़ो चर्च का जन्म 14 जून 1903 को वॉशिंगटन, डीसी में हुआ। आरंभिक
बाल्यावस्था में एयर गन से हुई एक दुर्घटना के कारण चर्च की एक आँख की दृष्टि
चली गई। उन्होंने 1920 में कनेक्टिकट से प्रारंभिक विद्यालय पूरा किया और उसी
वर्ष प्रिंसटन में विश्वविद्यालयी शिक्षा आरंभ की। 1927 में उनका डॉक्टरेट
अध्ययन पूरा हुआ। कुछ वर्ष विदेश में बिताकर चर्च प्रिंसटन लौटे। वे अत्यंत
विनम्र और सावधान माने जाते थे। श्यामपट्ट पर उनका लेखन निष्कलंक होता था और
महत्त्वपूर्ण कागज़ों को सुरक्षित रखने के लिए वे उन पर सावधानी से ड्यूको
सीमेंट (एक पारदर्शी गोंद) की परत चढ़ाते थे। अकादमिक कार्यों के बाहर उन्हें
विज्ञान-कथा पत्रिकाएँ पढ़ना पसंद था और लेखन में कोई त्रुटि दिखने पर वे
संपादकों को लिखने से नहीं हिचकते थे।

चर्च की अकादमिक उपलब्धियाँ महान थीं। अपने विद्यार्थियों स्टीफन क्लीन और
बार्कली रॉसर के साथ उन्होंने प्रभावी गणनीयता का सिद्धांत, लैम्ब्डा कलन,
ट्यूरिंग द्वारा ट्यूरिंग मशीन के विकास से स्वतंत्र रूप से विकसित किया।
गणनीयता की दोनों परिभाषाएँ तुल्य हैं और उनसे वह कथन उत्पन्न होता है जिसे अब
\emph{चर्च--ट्यूरिंग प्रस्थापना} कहते हैं: प्राकृतिक संख्याओं का कोई फलन
तभी और केवल तभी प्रभावी रूप से गणनीय है जब वह ट्यूरिंग मशीन (या लैम्ब्डा
कलन) द्वारा गणनीय हो। उन्होंने वह परिणाम भी सिद्ध किया जिसे अब \emph{चर्च
का प्रमेय} कहते हैं: प्रथम-क्रम सूत्रों की वैधता की निर्णय समस्या असमाधेय
है।

चर्च वृद्धावस्था तक अपना कार्य करते रहे। 1967 में वे प्रिंसटन छोड़कर यूसीएलए
गए, जहाँ 1990 में सेवानिवृत्ति तक प्राध्यापक रहे। चर्च का निधन 1 अगस्त 1995
को 92 वर्ष की आयु में हुआ।

\begin{reading} 
चर्च की संक्षिप्त जीवनी के लिए \citet{EndertonND} देखें। लैम्ब्डा कलन और
एंटशाइडुंग्सप्रॉब्लेम (चर्च की प्रस्थापना) पर चर्च के मूल लेख
\citet{Church1936,Church1936a} हैं। \citet{Aspray1984} में 1930 के दशक के
प्रिंसटन गणित समुदाय के बारे में चर्च का साक्षात्कार दर्ज है। चर्च ने 1936
से 1979 तक \emph{Journal of Symbolic Logic} के लिए पुस्तक समीक्षाओं की एक
शृंखला लिखी। वे सभी जॉन मैकफार्लेन की वेबसाइट पर अभिलेखित हैं
\citep{MacFarlane2015}।
\end{reading} 
\end{document}
